```latex
\documentclass[a4paper, 10pt]{article}
\usepackage[utf8]{inputenc}
\usepackage[T1]{fontenc}
\usepackage[margin=1in]{geometry}
\usepackage{enumitem}
\usepackage{hyperref}
\usepackage{helvet}
\renewcommand{\familydefault}{\sfdefault}

% Configuración de colores y enlaces
\hypersetup{
    colorlinks=true,
    linkcolor=black,
    citecolor=black,
    urlcolor=black
}

\begin{document}

% Encabezado
\begin{center}
    \textbf{\Large Sebastián Díaz de la Fuente}
    
    \vspace{0.2cm}
    {\small Ingeniero Civil en Informática | Especialista en Automatización y Arquitectura Backend}
    
    \vspace{0.4cm}
    {\footnotesize Foco en LLMs, Orquestación de Agentes y APIs RESTful}
\end{center}

\vspace{0.3cm}

% Resumen
\section*{Resumen Profesional}
Experto en diseño de arquitecturas backend escalables y automatización de flujos de trabajo complejos. Especialista en el desarrollo de APIs RESTful, ETL/ELT y despliegues de microservicios en Kubernetes. Cuenta con experiencia probada en el diseño de sistemas multi-agente utilizando LangGraph y Pydantic, aplicando principios de Clean Code para asegurar la robustez y validación estricta de datos. Busco aplicar mi experiencia en desarrollo backend y orquestación de sistemas para aportar valor a proyectos innovadores.

\vspace{0.3cm}

% Experiencia
\section*{Experiencia Profesional}

\textbf{Data Scientist \& Automation Analyst} \hfill \textit{Enero 2025 -- Presente}
\begin{itemize}[noitemsep, leftmargin=*]
    \item Diseño e implementación de \textbf{APIs RESTful} asíncronas para exponer modelos predictivos y facilitar la integración de sistemas.
    \item Desarrollo de orquestaciones \textbf{ETL/ELT} para la ingesta masiva de datos, optimizando los tiempos de procesamiento y limpieza.
    \item Creación de sistemas multi-agente persistentes utilizando \textbf{LangGraph} y \textbf{PostgreSQL}, automatizando procesos de auditoría forense.
    \item Implementación de validación de datos estricta con \textbf{Pydantic} (Structured Output) para reducir la tasa de error en modelos de OCR a menos del 1\%.
    \item Despliegue de infraestructura de inferencia contenida utilizando \textbf{Docker} y orquestación escalable para \textbf{Kubernetes}.
    \item Desarrollo de la herramienta CLI \textit{Headhunter}, un agente de ventas autónomo que automatiza la prospección.
\end{itemize}

\vspace{0.2cm}

\textbf{Data Analytics Intern} \hfill \textit{Junio 2024 -- Diciembre 2024}
\begin{itemize}[noitemsep, leftmargin=*]
    \item Automatización de la generación de inteligencia de negocio mediante consultas SQL avanzadas, optimizando reportes gerenciales en un 40\%.
    \item Desarrollo de scripts de automatización en entornos Linux y \textbf{Java Script} para optimizar flujos críticos de Recursos Humanos.
\end{itemize}

\vspace{0.2cm}

\textbf{Profesor y Ayudante} \hfill \textit{2020 -- 2024}
\begin{itemize}[noitemsep, leftmargin=*]
    \item Instrucción de programación y comunicación técnica, asegurando la comprensión de conceptos complejos por parte de estudiantes.
\end{itemize}

\textbf{Secretario Ejecutivo Federación Universitaria} \hfill \textit{2020 -- 2024}
\begin{itemize}[noitemsep, leftmargin=*]
    \item Gestión integral de iniciativas estudiantiles, enfocada en la retención y bienestar universitario.
\end{itemize}

\vspace{0.3cm}

% Habilidades
\section*{Habilidades Técnicas}
\begin{itemize}[noitemsep, leftmargin=*]
    \item \textbf{Backend & Lógica:} Python, \textbf{Java Script}, SQL, RESTful APIs, ETL/ELT, CLI Development, Pydantic.
    \item \textbf{Infraestructura:} \textbf{Docker}, \textbf{Kubernetes}, PostgreSQL.
    \item \textbf{Inteligencia Artificial:} LangGraph, LLMs, RAGs, Automation, OCR.
    \item \textbf{Metodologías:} Agile, DevOps, Buenas Prácticas de Código (Clean Code).
\end{itemize}

\vspace{0.3cm}

% Educación
\section*{Educación}
\begin{itemize}[noitemsep, leftmargin=*]
    \item \textbf{Ingeniería Civil en Informática} | Universidad Adolfo Ibáñez (UAI) | 2020 -- 2024
\end{itemize}

\vspace{0.3cm}

% Idiomas
\section*{Idiomas}
\begin{itemize}[noitemsep, leftmargin=*]
    \item Español (Nativo)
\end{itemize}

\end{document}
```